\documentclass[11pt,a4paper]{article}
\usepackage{amsmath,amssymb}
\usepackage{booktabs}
\usepackage{microtype}
\usepackage{parskip}

\title{VDW14 Tick-Counter Experiment\\
  \large CDBoost Adaptation --- Portless PDEVS}
\author{Damián Vicino}
\date{2026}

\begin{document}
\maketitle

\section*{Reference}
Vicino, Dalle, Wainer.
\textit{A Data Type for Discretized Time Representation in DEVS.}
SIMUTOOLS 2014. \texttt{hal-01055555}.

\section*{Canonical Model}

The tick-counter coupled model $C$ demonstrates causality errors from
floating-point time representation.

\[
  C = \langle X_C,\, Y_C,\, D,\, \{M_d\},\, \text{EIC},\, \text{EOC},\, \text{IC} \rangle
\]
\[
  D = \{G_{1/10},\; G_{1},\; K\},\quad
  X_C = \emptyset,\quad
  Y_C \subseteq \mathbb{N}
\]
\[
  \text{IC} = \{(G_{1/10}.\texttt{out} \to K.\texttt{add}),\;
               (G_{1}.\texttt{out} \to K.\texttt{reset})\}
\]
\[
  \text{EOC} = \{(K.\texttt{sum} \to C.\texttt{out})\}
\]

Generator $G_p = \langle S_G, ta_G, \delta^G_{\text{int}}, \lambda_G \rangle$:
\[
  S_G = \mathbb{R}^+,\quad
  ta_G(s) = p,\quad
  \delta^G_{\text{int}}(s) = s,\quad
  \lambda_G(s) = \mathit{tick}
\]

Counter $K = \langle S_K, ta_K, \delta^K_{\text{int}},
\delta^K_{\text{ext}}, \delta^K_{\text{con}}, \lambda_K \rangle$,
$S_K = \mathbb{N} \times \{\text{idle},\, \text{ready}\}$:
\begin{align*}
  ta_K(n,\text{idle})  &= \infty \qquad ta_K(n,\text{ready}) = 0 \\
  \delta^K_{\text{ext}}(s,e,x) &=
    \begin{cases}
      (n + 1,\; \text{idle})  & x \in \texttt{add}\\
      (n,\;     \text{ready}) & x \in \texttt{reset}
    \end{cases} \\
  \delta^K_{\text{int}}(n,\text{ready}) &= (0,\, \text{idle}) \\
  \lambda_K(n, \text{ready})            &= n
\end{align*}

Expected invariant: every output of $K$ equals $10$.

\section*{CDBoost Adaptation}

CDBoost implements a \emph{portless} PDEVS: all submodels in a coupled model
share a single message type \texttt{MSG}.  Routing is expressed as a list of
\texttt{(sender, receiver)} pointer pairs; there are no named ports.

Because $K$ receives two semantically distinct signals (\texttt{add} and
\texttt{reset}) on a single channel, the message \emph{value} must distinguish
them.  The adopted convention is:

\begin{center}
\begin{tabular}{@{}ll@{}}
  \toprule
  \texttt{MSG} value & Semantics \\
  \midrule
  $v \neq 0$ & add $v$ to the counter \\
  $0$        & output current count, then reset to zero \\
  \bottomrule
\end{tabular}
\end{center}

This matches the \texttt{infinite\_counter} model shipped with CDBoost.
$G_{1/10}$ outputs $1$; $G_{1}$ outputs $0$.  Both are routed to $K$ via IC.

\paragraph{Confluence.}
When $G_{1/10}$ and $G_1$ are simultaneously imminent---which occurs at every
integer second under exact arithmetic---the confluence function determines
whether the add or the reset is processed first:
\[
  \delta^K_{\text{con}}(s, e, x) = \delta^K_{\text{ext}}
      \!\left(\delta^K_{\text{int}}(s),\; 0,\; x\right)
\]
This rule applies the internal transition first (clear the ready flag if set),
then applies the external input, which corresponds to processing \texttt{add}
before \texttt{reset} when $G_{1/10} \prec G_1$.

\paragraph{Time representation.}
CDBoost requires a \texttt{TIME} type with a static \texttt{Inf()} method;
primitive \texttt{float} and \texttt{boost::rational} do not satisfy this
directly and need thin wrappers (\texttt{cdboost::float\_time},
\texttt{cdboost::rational\_time}).

\section*{Expected Observations}

\begin{itemize}
  \item \textbf{Rational time:} $K$ always outputs $10$.  The histogram has a
        single bar.
  \item \textbf{Float time:} $G_{1/10}$'s period $0.1$ is not exactly
        representable in IEEE~754.  Accumulated rounding error occasionally
        shifts a $G_{1/10}$ tick across a 1-second boundary; $K$ then outputs
        $9$ or $11$.  The effect is stochastic and grows with simulation length.
\end{itemize}

\end{document}
